\documentclass[../main.tex]{subfiles}
\usepackage[utf8]{inputenc}
\usepackage[T1]{fontenc}
\usepackage[indonesian]{babel}

\begin{document}

\chapter{Analisis Algoritma Block Cipher: AES}

\begin{learningoutcome}
Setelah mempelajari bab ini, mahasiswa diharapkan mampu:
\begin{itemize}
    \item Menjelaskan latar belakang pengembangan \textit{Advanced Encryption Standard} (AES) sebagai pengganti DES (Sub-CPMK 1.2).
    \item Menganalisis spesifikasi teknis AES, termasuk varian panjang kunci (128, 192, dan 256 bit) (Sub-CPMK 1.2).
    \item Menguraikan mekanisme kerja internal AES, termasuk transformasi \textit{SubBytes}, \textit{ShiftRows}, \textit{MixColumns}, dan \textit{AddRoundKey} (Sub-CPMK 1.2).
    \item Memahami konsep \textit{state} dan representasi data dalam matriks $4 \times 4$ pada AES (Sub-CPMK 1.2).
    \item Menjelaskan proses pembangkitan kunci putaran (\textit{Key Expansion}) pada AES (Sub-CPMK 1.2).
    \item Memahami dasar matematika $\text{GF}(2^8)$ yang digunakan dalam operasi AES (Sub-CPMK 1.2).
\end{itemize}
\end{learningoutcome}

\section{Latar Belakang AES}
Kelemahan utama DES terletak pada panjang kuncinya (56 bit) yang terlalu pendek untuk menghadapi serangan \textit{brute force} modern. Hal ini mendorong \textit{National Institute of Standards and Technology} (NIST) untuk mengadakan kompetisi terbuka guna mencari standar enkripsi blok baru yang lebih kuat, efisien, dan publik.

Pada tahun 2001, algoritma \textbf{Rijndael}, yang dirancang oleh Vincent Rijmen dan Joan Daemen, dipilih dan ditetapkan sebagai \textbf{Advanced Encryption Standard (AES)}.

\subsection{Persyaratan dan Spesifikasi AES}
AES dirancang dengan beberapa kriteria utama:
\begin{itemize}
    \item \textbf{Sifat Algoritma}: Harus merupakan \textit{symmetric block cipher} dengan rancangan yang sepenuhnya publik.
    \item \textbf{Ukuran Blok}: Tetap pada 128 bit.
    \item \textbf{Panjang Kunci}: Fleksibel, mendukung 128, 192, dan 256 bit.
    \item \textbf{Implementasi}: Harus efisien baik dalam perangkat lunak (\textit{software}) maupun perangkat keras (\textit{hardware}).
\end{itemize}

Varian AES ditentukan oleh panjang kuncinya:
\begin{itemize}
    \item \textbf{AES-128}: Kunci 128 bit, 10 putaran.
    \item \textbf{AES-192}: Kunci 192 bit, 12 putaran.
    \item \textbf{AES-256}: Kunci 256 bit, 14 putaran.
\end{itemize}

\section{Mekanisme Kerja AES}
Berbeda dengan DES yang menggunakan Jaringan Feistel, AES menggunakan struktur \textbf{Substitusi-Permutasi} (\textit{Substitution-Permutation Network}). Data diproses dalam orientasi \textit{byte} menggunakan matriks dua dimensi berukuran $4 \times 4$ yang disebut \textbf{state}.

\subsection{Representasi State}
Blok plainteks 128-bit dipetakan ke dalam matriks \texttt{state} $4 \times 4$ sebagai berikut:
\[
\text{State} = \begin{bmatrix}
s_{0,0} & s_{0,1} & s_{0,2} & s_{0,3} \\
s_{1,0} & s_{1,1} & s_{1,2} & s_{1,3} \\
s_{2,0} & s_{2,1} & s_{2,2} & s_{2,3} \\
s_{3,0} & s_{3,1} & s_{3,2} & s_{3,3}
\end{bmatrix}
\]

\subsection{Tahapan Enkripsi}
Proses enkripsi AES terdiri dari putaran awal, beberapa putaran utama, dan satu putaran akhir.

\subsubsection{Initial Round}
\begin{itemize}
    \item \textbf{AddRoundKey}: Melakukan operasi XOR antara \texttt{state} awal (plainteks) dengan kunci awal (\textit{cipher key}).
\end{itemize}

\subsubsection{Main Rounds (Putaran Utama)}
Setiap putaran utama terdiri dari empat transformasi:
\begin{enumerate}
    \item \textbf{SubBytes}: Substitusi setiap byte dalam \texttt{state} menggunakan tabel \textbf{S-box} non-linear. Tahap ini memberikan efek \textit{confusion}.
    \item \textbf{ShiftRows}: Pergeseran siklik (\textit{wrapping}) pada baris-baris matriks \texttt{state}. Baris 0 tetap, baris 1 geser 1 byte, baris 2 geser 2 byte, dan baris 3 geser 3 byte. Tahap ini memberikan efek \textit{diffusion}.
    \item \textbf{MixColumns}: Operasi perkalian matriks antara \texttt{state} dengan matriks konstanta tetap dalam field $\text{GF}(2^8)$. Tahap ini semakin memperkuat difusi antar kolom.
    \item \textbf{AddRoundKey}: XOR antara \texttt{state} hasil transformasi dengan kunci putaran (\textit{round key}) saat itu.
\end{enumerate}

\subsubsection{Final Round (Putaran Akhir)}
Putaran terakhir hampir sama dengan putaran utama, namun \textbf{tanpa tahap MixColumns}:
\[ \text{SubBytes} \rightarrow \text{ShiftRows} \rightarrow \text{AddRoundKey} \]

\section{Pembangkitan Kunci Putaran (Key Expansion)}
Fungsi \textit{KeyExpansion()} membangkitkan seluruh kunci putaran yang diperlukan dari satu kunci eksternal. Proses ini melibatkan operasi:
\begin{itemize}
    \item \textbf{RotWord}: Pergeseran sirkuler byte dalam sebuah \textit{word}.
    \item \textbf{SubWord}: Substitusi byte menggunakan S-box.
    \item \textbf{Rcon (Round Constant)}: XOR dengan konstanta putaran untuk mencegah simetri antar kunci putaran.
\end{itemize}

\section{Dasar Matematika: Galois Field $\text{GF}(2^8)$}
Operasi dalam AES, terutama \textit{MixColumns} dan S-box, dilakukan di dalam \textbf{Galois Field $\text{GF}(2^8)$}, yang merupakan medan berhingga biner.

\subsection{Representasi Polinomial}
Setiap byte (8 bit) direpresentasikan sebagai polinomial derajat 7 dengan koefisien dalam $\text{GF}(2)$:
\[ a = b_7x^7 + b_6x^6 + \dots + b_1x + b_0 \]

\subsection{Operasi Aritmetika}
\begin{itemize}
    \item \textbf{Penjumlahan}: Dilakukan dengan operasi XOR bitwise.
    \item \textbf{Perkalian}: Perkalian polinomial yang hasilnya kemudian direduksi menggunakan \textit{irreducible polynomial} $f(x) = x^8 + x^4 + x^3 + x + 1$.
\end{itemize}

\begin{obeactivity}
\textbf{Aktivitas 8.1: Simulasi Transformasi AES}
1. Ambil sebuah blok data 16-byte (misal: "Kriptografi-AES!").
2. Representasikan data tersebut ke dalam matriks \texttt{state} $4 \times 4$.
3. Lakukan simulasi tahap \textit{ShiftRows} secara manual pada matriks tersebut.
4. Amati bagaimana posisi byte berubah dan diskusikan pengaruhnya terhadap difusi data.
\end{obeactivity}

\begin{obereflection}
\textbf{Latihan 8.1}
1. Mengapa tahap \textit{MixColumns} dihilangkan pada putaran terakhir AES?
2. Jelaskan bagaimana kombinasi \textit{SubBytes}, \textit{ShiftRows}, dan \textit{MixColumns} secara kolektif mengimplementasikan prinsip \textit{confusion} dan \textit{diffusion} Shannon.
3. Berikan analisis singkat: mengapa AES-256 lebih tahan terhadap serangan \textit{quantum computing} (misal: algoritma Grover) dibandingkan AES-128?
\end{obereflection}

\begin{obeassessment}
\textbf{Kuis Bab 8}
1. Apa perbedaan utama antara struktur internal DES (Feistel) dan AES (Substitusi-Permutasi)?
2. Jelaskan peran \textit{S-box} dalam AES dan bagaimana ia dibentuk secara matematis.
3. Mengapa operasi dalam AES dilakukan dalam $\text{GF}(2^8)$ dan bukan aritmetika bilangan bulat biasa?
\end{obeassessment}

\begin{competencychecklist}
\begin{tabular}{|l|c|}
\hline
Indikator Kompetensi & Tercapai \\
\hline
Saya dapat menjelaskan alasan penggantian DES menjadi AES & [ ] \\
\hline
Saya dapat membedakan varian AES-128, 192, dan 256 & [ ] \\
\hline
Saya dapat menguraikan alur kerja \textit{SubBytes, ShiftRows, MixColumns} & [ ] \\
\hline
Saya dapat menjelaskan konsep matriks \textit{state} pada AES & [ ] \\
\hline
Saya dapat memahami dasar operasi aritmetika pada $\text{GF}(2^8)$ & [ ] \\
\hline
\end{tabular}
\end{center}
\end{competencychecklist}

\end{document}
